\documentclass[12pt, letterpaper]{article}

%-------Packages---------
\usepackage{amssymb,amsfonts}
\usepackage{mathptmx}
\usepackage{amsthm}
\usepackage{graphicx}
\usepackage[all,arc]{xy}
\usepackage[colorlinks=true, allcolors=blue]{hyperref}
\usepackage{enumerate}
\usepackage{bbm}
\usepackage{dsfont}
\usepackage{mathrsfs}
\usepackage{tikz-cd}
\usepackage{setspace}
\usepackage{import}
\usepackage[utf8]{inputenc}
\usepackage{hyperref}
\usepackage{tikz}
\usepackage{float}
\usepackage{mdframed}
\usepackage[nocheck]{fancyhdr}
\usepackage[nointegrals]{wasysym}

\usepackage[left=1in,top=1in,right=1in,bottom=1in]{geometry}

\usepackage{mathtools}
\usepackage{setspace}
\usepackage{cleveref}
\usepackage{multicol}

%--------Theorem Environments--------
%theoremstyle{plain} --- default
\newmdtheoremenv{theorem}{Theorem}[subsection]
\newmdtheoremenv{corollary}[theorem]{Corollary}
\newtheorem{proposition}[theorem]{Proposition}
\newmdtheoremenv{lemma}[theorem]{Lemma}
\newtheorem{conj}[theorem]{Conjecture}
\newtheorem{quest}[theorem]{Question}

\theoremstyle{definition}
\newmdtheoremenv{definition}[theorem]{Definition}
\newmdtheoremenv{definitions}[theorem]{Definitions}
\newtheorem{example}[theorem]{Example}
\newtheorem{algorithm}[theorem]{Algorithm}
\newtheorem{rmk}[theorem]{Remark}
\newtheorem{exercise}[theorem]{Exercise}

%----------- my commands -------------
\newcommand{\C}{\mathbb{C}}
\newcommand{\R}{\mathbb{R}}
\newcommand{\Q}{\mathbb{Q}}
\newcommand{\Z}{\mathbb{Z}}
\newcommand{\N}{\mathbb{N}}
\renewcommand{\P}{\mathbb{P}}
\newcommand{\contr}{\Rightarrow \Leftarrow}
\newcommand{\HRule}[1]{\rule{\linewidth}{#1}}
\newcommand{\s}{\(\,\)}
\renewcommand{\t}[1]{\text{#1}}
\newcommand{\ang}[1]{\left\langle #1 \right\rangle}

% Meta data
\setstretch{1.2}

\begin{document}

\pagestyle{fancy}
\fancyhf{}
\fancyhead[LOE]{\slshape{\textbf{Vincent Ferrara}}}
\fancyhead[ROE]{\thepage}

\subsection*{Problem 1}
\large$D_4$
\[
\begin{array}{c|ccccc}
\text{Class} & \{e\} & \{r^2\} & \{r, r^3\} & \{s, sr^2\} & \{sr, sr^3\} \\
\hline
\chi_1 & 1 & 1 & 1 & 1 & 1 \\
\chi_2 & 1 & 1 & -1 & 1 & -1 \\
\chi_3 & 1 & 1 & 1 & -1 & -1 \\
\chi_4 & 1 & 1 & -1 & -1 & 1 \\
\chi_5 & 2 & -2 & 0 & 0 & 0 \\
\end{array}
\]
\large$Q$
\[
\begin{array}{c|ccccc}
\text{Class} & \{1\} & \{-1\} & \{i, -i\} & \{j, -j\} & \{k, -k\} \\
\hline
\chi_1 & 1 & 1 & 1 & 1 & 1 \\
\chi_2 & 1 & 1 & -1 & 1 & -1 \\
\chi_3 & 1 & 1 & 1 & -1 & -1 \\
\chi_4 & 1 & 1 & -1 & -1 & 1 \\
\chi_5 & 2 & -2 & 0 & 0 & 0 \\
\end{array}
\]
They are the same table, but they are not isomorphic.

\newpage
\subsection*{Problem 2}
If $n$ is even then from HW 3 it has classes $\{\{e\},\{x,x^{-1}\},\{x^2,x^{-2}\},\dots,\{y,yx^2,\dots\},\{yx,yx^3,\dots\}\}$\\
Thus $D_n$ has $n/2+3$ classes thus we have $n/2+3$ representations.\\
Consider the 1-d case thus $\rho: G \rightarrow \C^\times$\\
$\rho(y)=\{-1,1\}$, consider $\rho(yxy)=\rho(y)^2\rho(x)=\rho(x)$ thus $\rho(r)^2=1$ also thus $\rho(r)=\{-1,1\}$\\
Thus we have 4 1-d representations. Trivial, $\{\rho(x)=-1,\rho(y)=1$\}, $\{\rho(x)=1, \rho(y)=-1\}$, $\{\rho(x)=-1, \rho(y)=-1\}$\\
2-d representations. Let $\omega = e^{2i\pi/n}$\\
Define them as:
\[
\rho^h(x)=\begin{pmatrix}
    \omega^h & 0\\
    0 & \omega^{-h}
\end{pmatrix}
\quad
\rho^h(y)=\begin{pmatrix}
    0 & 1\\
    1 & 0
\end{pmatrix}
\]
Consider $T=\begin{pmatrix}
    0 & 1\\
    1 & 0
\end{pmatrix}$ $Tp^h(x)=p^{n-h}(x)T$ and $Tp^h(y)=p^{n-h}(y)T$ Thus by Schur's Lemma we know that $p^h \cong p^{n-h}$\\
Thus consider $0 \leq h \leq n/2$\\
For $h=0$ consider $\ang{\chi,\chi}=\dfrac{1}{2n}\sum_{g\in D_n}|\chi(g)|^2=2$ thus this is not irreducible\\
For $h=n/2$ consider $\ang{\chi,\chi}=\dfrac{1}{2n}\sum_{g\in D_n}|\chi(g)|^2=2$ thus this is not irreducible.\\
For $0 < h < n/2$. Since $\omega^h \neq \omega^{-h}$ the only lines of $C^2$ which are G-invariant under $\rho^h(y)$ are $y=0$ and $x=0$,
but these lines are not G-invariant under $\rho^h(x)$ thus these are irreducible.\\
Thus, we have $4+(n/2-1)=(n/2)+3$ representations thus we are done.\\
\phantom{text}\\
If $n$ is odd then we have $(n+1)/2+1$\\
Consider the 1-d cases\\
We have the same cases as above where $\rho(y)=\{-1,1\}$ and $\rho(x)=\{-1,1\}$ however since $(-1)^n=-1$ thus this case does not work thus we only have two 1-d representations\\
Trivial and $\rho(y)=-1, \rho(x)=1$\\
Consider the same matrices\\
\[
\rho^h(x)=\begin{pmatrix}
    \omega^h & 0\\
    0 & \omega^{-h}
\end{pmatrix}
\quad
\rho^h(y)=\begin{pmatrix}
    0 & 1\\
    1 & 0
\end{pmatrix}
\]
Again $p^h \cong p^{n-h}$\\
Thus consider $0 \leq h \leq (n-1)/2$\\
For $h=0$ consider $\ang{\chi,\chi}=\dfrac{1}{2n}\sum_{g\in D_n}|\chi(g)|^2=2$ thus again this is irreducible\\
For $0 < h \leq (n-1)/2$. Again since $\omega^h \neq \omega^{-h}$ the only lines of $C^2$ which are G-invariant under $\rho^h(y)$ are $y=0$ and $x=0$,
but these lines are not G-invariant under $\rho^h(x)$ thus these are irreducible.\\
Thus, we have $2+(n-1)/2=(n+1)/2+1$ thus we are done.

\newpage
\subsection*{Problem 3}
\large$S_4$
\[
\begin{array}{c|ccccc}
\text{Class} & \{e\} & \{(12)\} & \{(12)(34)\} & \{(123)\} & \{(1234)\} \\
\hline
\chi_{triv} & 1 & 1 & 1 & 1 & 1 \\
\chi_{sgn} & 1 & -1 & 1 & 1 & -1 \\
\chi_3 & 2 & 0 & 2 & -1 & 0 \\
\chi_{std} & 3 & 1 & -1 & 0 & -1 \\
\chi_5 & 3 & -1 & -1 & 0 & -1 \\
\end{array}
\]
$\chi_1$:trivial\\
$\chi_2$:sign\\
Consider the representation of $S_4$ on $\C^4$ where you just permute the elements.
Thus, we can break it into the trivial representation plus its orthogonal complement well $\ang{u,1}=u_1+u_2+u_3+u_4=0$\\
Thus the orthogonal complement is the standard representation of $S_4$ on $\C^3$\\
Assume there is a nontrivial invariant subspace of this called $W$.\\
If dim$(W)=1$\\
Then it is spanned by a single vector $w=(w_1,w_2,w_3,w_4)$, and also we have that $w_1+w_2+w_3+w_4=0$.
However, $S_4$ permutes the coordinates of $w$ thus for $W$ to be invariant this implies $w_1=w_2=w_3=w_4$; however this is not possible since $w_1+w_2+w_3+w_4=0$\\
If dim$(W)=2$\\
Then it is spanned by two vector $v,w$. However, this is not possible since if we permute the elements of things on this plane $W$ we can easily get stuff off the plane for example for $w$ just swap any two of the elements, and we will get something off the plane.\\
Thus, this is $G$-invariant.\\
Thus, we can get the character by this $\chi = \chi_{nat}-\chi_{triv}$\\
$\chi(e)=\chi_{nat}(e)-\chi_{triv}(e)=4-1=3$\\
$\chi((12))=\chi_{nat}((12))-\chi_{triv}((12))=2-1=1$\\
$\chi((123))=\chi_{nat}((123))-\chi_{triv}((123))=1-1=0$\\
$\chi((1234))=\chi_{nat}((1234))-\chi_{triv}((1234))=0-1=-1$\\
$\chi((12)(34))=\chi_{nat}((12)(34))-\chi_{triv}((12)(34))=0-1=-1$\\
Now using the hint multiply by sign representation, and we get $\chi(g)=\chi_{std}(g)\chi_{sgn}(g)$\\
And to confirm it is irreducible $\ang{\chi,\chi}=\dfrac{1}{24}(9+6+0+6+3)=1$\\
Now since I have all but one row left I can get the remaining row which is 2-d\\
$\chi_{triv}\downarrow_{A_4}=\chi_{triv}'$\\
$\chi_{sgn}\downarrow_{A_4}=\chi_{triv}'$\\
$\chi_{std}\downarrow_{A_4}=\chi_{5}\downarrow_{A_4}=\chi_{std}'$\\
$\chi_{std}'(e)=3$\\
$\chi_{std}'((12)(34))=-1$\\
$\chi_{std}'((123))=0$\\
$\chi_{std}'((132))=0$\\
Consider $\ang{\chi_{std}',\chi_{std}'}=1$\\
$\chi_{3}\downarrow_{A_4}=\chi_{V}'$\\
$\chi_{V}'(e)=2$\\
$\chi_{V}'((12)(34))=2$\\
$\chi_{V}'((123))=-1$\\
$\chi_{V}'((132))=-1$\\
Well I have 2/4 characters of $A_4$ and the remaining must be 1-d since $1+a^2+b^2+3^2=12 \implies a^2+b^2=2$ thus $a=b=1$\\
Now we can solve for the last 3 spots in each row, and we get that $1,1,\omega,\omega^2$ call it $\chi_2'$ and $1,1,\omega^2,\omega$ call it $\chi_3'$ where it goes $e,(12)(34),(123),(132)$\\
Thus we can easily see that $\chi_{V}'=\chi_{2}'+\chi_{3}'$

\newpage
\subsection*{Problem 4}
\begin{enumerate}[(a)]
    \item Let $g$ be the generator of $C_4$ then let $R(g)=\begin{pmatrix}
        0 & 0 & 0 & 1\\
        1 & 0 & 0 & 0\\
        0 & 1 & 0 & 0\\
        0 & 0 & 1 & 0
    \end{pmatrix}$\\
    char($R(g)$)=$\lambda^4-1$\\
    Thus the eigenvalues of this are the 4th roots of unity.\\
    To decompose the regular representation into irreducible \( G \)-invariant subspaces using eigenvalues, let's proceed with the matrix representation of \( R(g) \) and analyze its eigenspaces.

The eigenvalues are the 4th roots of unity:
\[
\lambda_0 = 1, \quad \lambda_1 = i, \quad \lambda_2 = -1, \quad \lambda_3 = -i.
\]

Each eigenvalue corresponds to a 1-dimensional eigenspace, and these eigenspaces form \( G \)-invariant subspaces.
For \( \lambda_0 = 1 \):
The eigenspace is determined by solving \( (R(g) - I)v = 0 \). Explicitly:
\[
\begin{pmatrix}
-1 & 0 & 0 & 1 \\
1 & -1 & 0 & 0 \\
0 & 1 & -1 & 0 \\
0 & 0 & 1 & -1
\end{pmatrix}
\begin{pmatrix}
v_1 \\ v_2 \\ v_3 \\ v_4
\end{pmatrix}
= 0.
\]
The solution is \( v = (1, 1, 1, 1) \). This corresponds to the trivial representation \( \rho_0 \), where all elements of \( G \) act as 1.\\
For \( \lambda_1 = i \):
The eigenspace is determined by solving \( (R(g) - iI)v = 0 \). Explicitly:
\[
\begin{pmatrix}
-i & 0 & 0 & 1 \\
1 & -i & 0 & 0 \\
0 & 1 & -i & 0 \\
0 & 0 & 1 & -i
\end{pmatrix}
\begin{pmatrix}
v_1 \\ v_2 \\ v_3 \\ v_4
\end{pmatrix}
= 0.
\]
The solution gives a vector \( v = (-i, -1, i, 1) \). This corresponds to the irreducible representation \( \rho_{i} \), where \( g \) acts as multiplication by \( i \).\\
 For \( \lambda_2 = -1 \):
The eigenspace is determined by solving \( (R(g) + I)v = 0 \). Explicitly:
\[
\begin{pmatrix}
1 & 0 & 0 & 1 \\
1 & 1 & 0 & 0 \\
0 & 1 & 1 & 0 \\
0 & 0 & 1 & 1
\end{pmatrix}
\begin{pmatrix}
v_1 \\ v_2 \\ v_3 \\ v_4
\end{pmatrix}
= 0.
\]
The solution gives a vector \( v = (-1, 1, -1, 1) \). This corresponds to the irreducible representation \( \rho_{-1} \), where \( g \) acts as multiplication by \( -1 \).\\
 For \( \lambda_3 = -i \):
The eigenspace is determined by solving \( (R(g) - (-i)I)v = 0 \). Explicitly:
\[
\begin{pmatrix}
i & 0 & 0 & 1 \\
1 & i & 0 & 0 \\
0 & 1 & i & 0 \\
0 & 0 & 1 & i
\end{pmatrix}
\begin{pmatrix}
v_1 \\ v_2 \\ v_3 \\ v_4
\end{pmatrix}
= 0.
\]
The solution gives a vector \( v = (i, -1, -i, 1) \). This corresponds to the irreducible representation \( \rho_{-i} \), where \( g \) acts as multiplication by \( -i \).\\
Thus, the $R = \rho_{triv} \oplus \rho_i \oplus \rho_{-1} \oplus \rho_{-i}$

\item Let $a,b$ be the generators for $G$ then let
\[
R(a)=\begin{pmatrix}
    0 & 1 & 0 & 0\\
    1 & 0 & 0 & 0\\
    0 & 0 & 0 & 1\\
    0 & 0 & 1 & 0
\end{pmatrix}
\quad R(b)=\begin{pmatrix}
    0 & 0 & 1 & 0\\
    0 & 0 & 0 & 1\\
    1 & 0 & 0 & 0\\
    0 & 1 & 0 & 0
\end{pmatrix}
\]
Using the same methods we get that the eigenvalues are $-1,1$ for both and the eigenvectors for $R(a)$ are $v=(a,a,b,b)$, and $w=(-a,a,-b,b)$
, and eigenvectors for $R(b)$ are $v=(a,b,a,b)$, and $w=(-a,-b,a,b)$\\
Thus we can split it into 4 cases\\
\begin{itemize}
    \item Eigenvector $v=(1,1,1,1)$ this is the trivial space.
    \item Eigenvector $v=(1,-1,1,-1)$ this is the subspace that sends $\rho_1(a)=-1$ and $\rho_1(b)=1$
    \item Eigenvector $v=(1,1,-1,-1)$ this is the subspace that sends $\rho_2(a)=1$ and $\rho_2(b)=-1$
    \item Eigenvector $v=(1,-1,-1,1)$ this is the subspace that sends $\rho_3(a)=-1$ and $\rho_3(b)=-1$
\end{itemize}
Thus we get that $R=\rho_{triv} \oplus \rho_1 \oplus \rho_2 \oplus \rho_3$
\end{enumerate}
Thus I have verified the theorem.
\newpage
\subsection*{Problem 5}
Consider $\chi$ a nontrivial irreducible character on the table, and let $\rho$ be its representation.\\
Let ker$(\chi)=\{g \in G \mid \chi(\rho(g))=\chi(\rho(e))\}$\\
Let $x \in \t{ker}(\rho)$ thus $\rho(x)=id$ thus $\chi(\rho(x))=\chi(id)=\chi(\rho(e))$ thus $x \in \t{ker}(\chi)$\\
Let $x \in \t{ker}(\chi)$ then we know that the sum of $\rho(x)$ eigenvalues sum to the dimension of the space. From class, we know that the eigenvalues are all roots of unity thus $|\lambda|<=1$ thus to sum dimension of the space
amount of eigenvalues and get the dimension is only possible if all eigenvalues are 1 thus $\rho(x)=id$ thus $x \in \t{ker}(\rho)$\\
Thus ker$(\chi)=\t{ker}(\rho)$\\
Assume all irreducible characters of $G$ have trivial kernels.\\
Assume $G$ is not simple thus let $N$ be a non-trivial normal subgroup.
Consider the group $G\diagup N$ since this group is not trivial it will have one nontrivial character, thus we can compose this character with the natural map $\pi : G \rightarrow G\diagup N$ to get a character in $G$ thus we can see that $N$ is a subset of this character's kernel.
Then we can then reduce it to a sum of irreducible characters in $G$. If every irreducible character of $G$ has a trivial kernel, then the intersection of their kernels must also be trivial thus this is not possible since we know kernel of this composite character is not trivial,
thus at least one irreducible character of $G$ has a nontrivial kernel thus this is a contradiction. Thus, $G$ is simple.\\
Thus, from a table if you see any nontrivial irreducible characters with $\chi(g)=\chi(e)$ then you know $G$ is not simple, and if
it is simple then all nontrivial irreducible characters will have $\chi(g)\neq\chi(e)$ for all $g\in G$.

\newpage
\subsection*{Problem 6}
$\chi_A(t)=\det(tI-A)=(1-t)^2$
Consider the span of $v$ s.t. $v=(1,0)$\\
Thus, let $x \in \t{span}(v)$; thus, $x=cv$ for $c \in \mathbb{F}_p$.
Thus, consider $Ax=x$ since $A$'s only eigenvalue is 1, and thus span$(v) = V_1$ is a $G$-invariant subspace, and $V_1$ is irreducible since it is dim 1.\\
Thus, assume $V=V_1 \bigoplus W$ s.t. $W$ is another $G$-invariant subspace, and again it is irreducible since dim 1.
Thus, $W=\t{span}(w)$ for some $w \in V$ and again $Aw=w$ since its only eigenvalue is 1.
However, this would mean that $A$ is diagonalizable, but it is not thus $w,v$ cannot be lineally independent.\\
Thus, this is a contradiction and $V$ cannot be written as a direct sum of $G$-invariant subspaces.
\end{document}